% Generated by tools/build_new_dists_anthology_sections.py; do not edit by hand.

\section{Sporadic groups and represented covers}\label{proof:new-dist:sporadic}

\resultbox{\textbf{Canonical testing artifacts.}\par\smallskip Exact backend: \path{lambda_distributions/dists/sporadic_group_sigma_mgfs/verification.py}.\par Regression: \path{lambda_distributions/dists/sporadic_group_sigma_mgfs/test_verification.py}.\par Marimo: \path{marimo_notebooks/sporadic.py}.}

\subsection{A representation is part of the data}

A sporadic abstract group does not by itself determine a sigma--MGF.  Fix a
finite-dimensional complex representation
\[
\rho:G\longrightarrow GL(V),\qquad \chi(g)=\operatorname{Tr}\rho(g).
\]
For formal variables $\mathbf t=(t_1,t_2,\ldots)$ and power sums
$p_r=\sum_i t_i^r$, define
\begin{equation}\label{nd:sporadic:eq:def}
\MGF_{G,V}(\mathbf t)=\frac1{|G|}\sum_{g\in G}
\prod_{i\geq1}\det(I-t_i\rho(g))^{-1}.
\end{equation}
Different representations of the same group generally give different
functions.  This note groups together two representations arising from the
same $M_{11}$ action: the full permutation module and its deleted constituent.

\subsection{Universal class-power theorem}

Let $x_1,\ldots,x_s$ represent the conjugacy classes of $G$, and put
$z_a=|C_G(x_a)|$.  Then
\begin{equation}\label{nd:sporadic:eq:universal}
\boxed{
\MGF_{G,V}(\mathbf t)
=\sum_{a=1}^s\frac1{z_a}
\exp\!\left(\sum_{r\geq1}\frac{p_r}{r}\chi(x_a^r)\right).}
\end{equation}
For a partition $\tau=(\tau_1,\ldots,\tau_\ell)$, with
$m_\tau=t_1^{\tau_1}\cdots t_\ell^{\tau_\ell}$ before symmetrization,
\begin{equation}\label{nd:sporadic:eq:coefficient}
\boxed{[m_\tau]\MGF_{G,V}
=\dim\left(\bigotimes_{j=1}^{\ell}\Sym^{\tau_j}V\right)^G.}
\end{equation}
Thus class centralizers, class power maps, and one ordinary character suffice
to determine the complete sigma--MGF.

\subsubsection*{Proof}
If the eigenvalues of a matrix $A$ are $\lambda_1,\ldots,\lambda_n$, then,
as an identity of formal power series,
\begin{align*}
\det(I-tA)^{-1}
&=\prod_{j=1}^n(1-t\lambda_j)^{-1}\\
&=\exp\!\left(\sum_{r\geq1}\frac{t^r}{r}
\operatorname{Tr}(A^r)\right).
\end{align*}
Multiplication over $t_i$, followed by $A=\rho(g)$ and
$\operatorname{Tr}(\rho(g)^r)=\chi(g^r)$, proves the elementwise equivalence
of the determinant and exponential integrands.  Grouping the average by
classes uses $|x_a^G|/|G|=1/z_a$ and proves \eqref{nd:sporadic:eq:universal}.

The coefficient of $t^d$ in $\det(I-t\rho(g))^{-1}$ is the character of
$\Sym^dV$ at $g$.  Therefore the coefficient indexed by $\tau$ is the
average of the character of $\bigotimes_j\Sym^{\tau_j}V$.  This group average
is the trace, hence the rank, of the Reynolds projector onto the fixed
subspace, proving \eqref{nd:sporadic:eq:coefficient}.  The proof is exact in every degree.

\subsection{The \texorpdfstring{$M_{11}$}{M11} matrix group}

The online ATLAS gives the following standard generators in the natural
11-point action:
\begin{align*}
a&=(2,10)(4,11)(5,7)(8,9),\\
b&=(1,4,3,8)(2,5,6,9).
\end{align*}
The verifier closes these permutations under multiplication and inversion;
it obtains exactly $7{,}920$ elements.  With $P_ge_j=e_{g(j)}$, this constructs
all $11\times11$ zero-one matrices $P_g$ rather than merely storing abstract
class data.  The standard-generator fingerprint is
\[
o(a)=2,\qquad o(b)=4,\qquad o(ab)=11,\qquad
o(ababbabbb)=5.
\]

For the deleted representation, let
\[
W=\{x\in\mathbb C^{11}:x_1+\cdots+x_{11}=0\},\qquad
f_j=e_j-e_{11}\quad(1\leq j\leq10).
\]
The code constructs the restriction $D_g=P_g|_W$ in this basis.  Explicitly,
\begin{equation}\label{nd:sporadic:eq:deleted-matrix}
(D_g)_{ij}=\delta_{i,g(j)}-\delta_{i,g(11)}
\quad(1\leq i,j\leq10),
\end{equation}
where a delta is zero if its second index is $11$.  Every entry is integral.
All $7{,}920$ matrices are distinct in each representation, and sampled
products satisfy $P_gP_h=P_{gh}$ and $D_gD_h=D_{gh}$ exactly.

\subsection{Exact formulas for the two \texorpdfstring{$M_{11}$}{M11} representations}

If $\lambda=1^{c_1}2^{c_2}\cdots$ is a permutation cycle shape, set
\[
Z_\lambda(\mathbf t)=\prod_{i\geq1}\prod_{r\geq1}
(1-t_i^r)^{-c_r},\qquad
\widetilde Z_\lambda(\mathbf t)=\prod_{i\geq1}(1-t_i)Z_\lambda(\mathbf t).
\]
The second identity follows from
$\mathbb C^{11}\cong\mathbf1\oplus W$.  Direct enumeration gives the eight
aggregated spectra below; the two order-eight classes and two order-eleven
classes may be combined because their spectra in these representations agree.

\begin{center}
\begin{tabular}{@{}lrrl@{}}
\toprule
class row & size & centralizer & cycle shape\\
\midrule
$1A$ & 1 & 7,920 & $1^{11}$\\
$2A$ & 165 & 48 & $1^3 2^4$\\
$3A$ & 440 & 18 & $1^2 3^3$\\
$4A$ & 990 & 8 & $1^3 4^2$\\
$5A$ & 1,584 & 5 & $1\,5^2$\\
$6A$ & 1,320 & 6 & $2\,3\,6$\\
$8AB$ & 1,980 & $8$ in each class & $1\,2\,8$\\
$11AB$ & 1,440 & $11$ in each class & $11$\\
\bottomrule
\end{tabular}
\end{center}

Consequently the natural permutation formula is
\begin{equation}\label{nd:sporadic:eq:permformula}
\boxed{\begin{aligned}
\MGF_{M_{11},\mathbb C^{11}}=\frac1{7920}\big(&
Z_{1^{11}}+165Z_{1^3 2^4}+440Z_{1^2 3^3}
+990Z_{1^3 4^2}\\
&+1584Z_{1\,5^2}+1320Z_{2\,3\,6}
+1980Z_{1\,2\,8}+1440Z_{11}\big).
\end{aligned}}
\end{equation}
The deleted formula is the same weighted sum with every $Z$ replaced by
$\widetilde Z$:
\begin{equation}\label{nd:sporadic:eq:delformula}
\boxed{\MGF_{M_{11},W}
=\left(\prod_{i\geq1}(1-t_i)\right)
\MGF_{M_{11},\mathbb C^{11}}.}
\end{equation}
Equation \eqref{nd:sporadic:eq:delformula} is a relation between generalized Molien
integrands and hence between their averages; it does not assert a tensor
factorization of invariant spaces.

The same cycle data recover every character of a power:
\begin{equation}\label{nd:sporadic:eq:powertrace}
\chi_{\rm perm}(g^r)=\sum_{d\mid r}d\,c_d(g),\qquad
\chi_W(g^r)=\chi_{\rm perm}(g^r)-1.
\end{equation}
Substitution in \eqref{nd:sporadic:eq:universal} gives a third, power-character evaluation
of \eqref{nd:sporadic:eq:permformula}--\eqref{nd:sporadic:eq:delformula}.

\subsection{Verification strategy}

The program uses four routes with deliberately different inputs.
\begin{enumerate}[leftmargin=*,itemsep=3pt]
\item \textbf{Direct matrices.}  For every generated matrix $A$, traces of
$A^r$ are computed by integer matrix multiplication.  Complete symmetric
characters are recovered from Newton's exact recurrence
\[
h_0(A)=1,\qquad
h_d(A)=\frac1d\sum_{r=1}^d\operatorname{Tr}(A^r)h_{d-r}(A).
\]
The target coefficient is then the literal average
$|G|^{-1}\sum_g\prod_jh_{\tau_j}(\rho(g))$.
\item \textbf{Class determinants.}  Only the eight table rows are used to
expand $\prod_r(1-t^r)^{-c_r}$; the deleted factor $1-t$ is then applied.
\item \textbf{Power characters.}  Equation \eqref{nd:sporadic:eq:powertrace} and the
same Newton recurrence expand the exponential formula, without using the
determinant factorization in the implementation.
\item \textbf{Configuration orbits.}  For selected permutation coefficients,
the two generators traverse the finite set
\[
\prod_j\MSet_{\tau_j}(\{1,\ldots,11\}).
\]
The number of orbits must equal \eqref{nd:sporadic:eq:coefficient}.  This route uses
neither the class table nor matrix traces.
\end{enumerate}

Finally, at three nonzero variables the code evaluates every determinant
numerically and compares the full average with the closed class formula.
This catches mistakes that could be hidden by a short coefficient truncation.

\subsection{Results}

All 29 partitions of positive total degree at most six were tested in each
representation, giving 58 exact three-way comparisons.  Representative rows
are shown here.

\begin{center}
\begin{tabular}{@{}llrrrc@{}}
\toprule
representation & $\tau$ & matrix & determinant & power character & result\\
\midrule
permutation & $(1)$ & 1 & 1 & 1 & \textcolor{teal}{\textbf{PASS}}\\
permutation & $(2)$ & 2 & 2 & 2 & \textcolor{teal}{\textbf{PASS}}\\
permutation & $(2,1)$ & 4 & 4 & 4 & \textcolor{teal}{\textbf{PASS}}\\
permutation & $(2,2)$ & 9 & 9 & 9 & \textcolor{teal}{\textbf{PASS}}\\
permutation & $(1,1,1,1)$ & 15 & 15 & 15 & \textcolor{teal}{\textbf{PASS}}\\
permutation & $(6)$ & 13 & 13 & 13 & \textcolor{teal}{\textbf{PASS}}\\
\addlinespace
deleted & $(1)$ & 0 & 0 & 0 & \textcolor{teal}{\textbf{PASS}}\\
deleted & $(2)$ & 1 & 1 & 1 & \textcolor{teal}{\textbf{PASS}}\\
deleted & $(2,1)$ & 1 & 1 & 1 & \textcolor{teal}{\textbf{PASS}}\\
deleted & $(2,2)$ & 3 & 3 & 3 & \textcolor{teal}{\textbf{PASS}}\\
deleted & $(1,1,1,1)$ & 4 & 4 & 4 & \textcolor{teal}{\textbf{PASS}}\\
deleted & $(6)$ & 5 & 5 & 5 & \textcolor{teal}{\textbf{PASS}}\\
\bottomrule
\end{tabular}
\end{center}

The separate generator-orbit counts for
$\tau=(1),(2),(1,1),(3),(2,1),(2,2),(1,1,1,1)$ are respectively
\[
1,\ 2,\ 2,\ 3,\ 4,\ 9,\ 15,
\]
and agree with the class coefficients.  In particular, the last value is the
Bell number $B_4=15$, as expected from the 4-transitivity of this action.

At $(t_1,t_2,t_3)=(0.037,-0.051,0.083)$, the full determinant averages are
\begin{center}
\begin{tabular}{@{}lrrr@{}}
\toprule
representation & direct matrix average & class formula & absolute error\\
\midrule
permutation & 1.0866701548337285 & 1.0866701548337498 & $2.13\times10^{-14}$\\
deleted & 1.0085468522143926 & 1.0085468522144532 & $6.06\times10^{-14}$\\
\bottomrule
\end{tabular}
\end{center}
The discrepancies are at floating-point roundoff scale; all coefficient and
group-law checks are exact integers.

\subsection{Other sporadic constructions and scope}

\subsubsection{Other ordinary representations}
For any of the 26 sporadic simple groups, and for any chosen ordinary complex
character, equation \eqref{nd:sporadic:eq:universal} applies verbatim.  For a nontrivial
irreducible representation of a sporadic simple group, the kernel is normal
and therefore trivial.  This observation does not select a representation;
one must still state which character is intended.

\subsubsection{Central covers}
If $\widetilde G$ is a central cover and a central element $z$ acts on $V$ as
the scalar $\omega(z)$, it acts on
$\bigotimes_j\Sym^{\tau_j}V$ as $\omega(z)^{|\tau|}$.  Hence
\[
[m_\tau]\MGF_{\widetilde G,V}=0
\quad\text{unless}\quad \omega(z)^{|\tau|}=1
\quad\text{for every }z\in Z(\widetilde G).
\]
Projective representations of the simple quotient must therefore be treated
as honest representations of the appropriate cover.

\subsubsection{Lattice and conjugation actions}
For a lattice action, class eigenvalues or frame shapes may be substituted
directly into \eqref{nd:sporadic:eq:def}.  For conjugation on the basis $\mathbb C[G]$,
the fixed-point character is $\chi(g^r)=|C_G(g^r)|$; M\"obius inversion then
recovers the conjugation cycle counts.  These are separate representation
choices and are not silently combined with the $M_{11}$ action tested here.

\subsubsection{Asymptotics}
There is no intrinsic limiting sequence over the sporadic simple groups:
there are only 26.  For the one-variable coefficients of either fixed faithful
$M_{11}$ representation, ordinary rational-function asymptotics do apply.
The identity matrix supplies the unique pole of order $\dim V$ at $t=1$, up
to root-of-unity oscillations of lower polynomial order, so
\begin{align*}
[t^d]\MGF_{M_{11},\mathbb C^{11}}(t)
&=\frac1{7920}\binom{d+10}{10}+O(d^9),\\
[t^d]\MGF_{M_{11},W}(t)
&=\frac1{7920}\binom{d+9}{9}+O(d^8).
\end{align*}
The exact sequences are quasipolynomial because every eigenvalue is a root of
unity.

\subsection{Reproduction and provenance}

From the repository root, run
\begin{verbatim}
.venv/bin/python -m lambda_distributions.dists.sporadic_group_sigma_mgfs.verification
.venv/bin/python -m pytest -q \
  lambda_distributions/dists/sporadic_group_sigma_mgfs/test_verification.py
.venv/bin/marimo edit marimo_notebooks/sporadic.py
\end{verbatim}
The generator data, group order, standard-generator fingerprint, representation
degree, and class centralizers are recorded by the online
\href{https://brauer.maths.qmul.ac.uk/Atlas/v3/spor/M11/}{ATLAS page for
$M_{11}$}; the natural action and its cycle types are recorded on the
\href{https://brauer.maths.qmul.ac.uk/Atlas/v3/permrep/M11G1-p11B0}{ATLAS
11-point representation page}.  The program independently enumerates the
group and reproduces every aggregated class multiplicity used in
\eqref{nd:sporadic:eq:permformula}.

This is a standalone inclusion-ready proof.  It does not modify or rebuild the
combined proof PDF.
